
% 
% \section{Discussion}
% \label{sec:discussion}

% Several limitations of our work point to promising future research directions. \textbf{(i)} The baseline trajectories used in our dataset are locally accurate via ARKit–LiDAR alignment but are not metrology-grade globally; incorporating motion-capture or total-station measurements would further improve absolute accuracy. \textbf{(ii)} We primarily diagnose failure modes rather than propose algorithms; progress likely demands scale-invariant scene representations, tightly-coupled visual–inertial fusion, and robust, non-Gaussian back-ends for long-range loop closure under scale drift. \textbf{(iii)} Beyond single sessions, cross-session and multi-robot scale consistency remains largely open and is directly supported by our benchmark design.

% 
\section{Conclusion}
\label{sec:conclusion}

In this work, we examined the fundamental challenges of three types of scale inconsistencies in deep monocular visual \ac{SLAM}. With ScaleMaster, which exposes intra-session scale drift and inter-session scale ambiguity in large indoor environments, we coupled traditional trajectory evaluation with a direct map-to-map assessment (e.g., Chamfer distance and Drop Rate). Across multiple baselines (including DROID-SLAM, MASt3R-SLAM, and VGGT-SLAM), our quantitative and qualitative results reveal severe scale-related failures under ScaleMaster's more realistic conditions, despite good performance on existing benchmarks. We hope this dataset and the accompanying baseline evaluations provide a solid platform for measuring and improving scale-consistent, reliable SLAM.
